\section{The open directions, and my assessment}\label{sec:directions}

The workbench explored its problem along many routes, tested quickly with AI systems. Some produced theorems, and some mapped out what a route needs. This section states, for each open direction, what is established and my assessment. The assessments are views held at this stage, with their reasons.

\paragraph{Locally gapped, globally gapless sequences.}
The spatial line (129 pages) builds sequences in which the box, the spacing and the coupling vary together, driven by the $S^6$ period data through the magnetic background of Section~\ref{sec:side}. Its conclusion (§28) states three results: along a fixed-coupling diagonal with $g_0^4\ge12288$ the magnetic excitation probability eventually has zero weight on every bounded energy interval; for fixed Wilson-loop states there is a positive high-energy mass and a failure of strong continuity; and on a selected weak-coupling diagonal the finite gaps close while the corrected magnetic probability escapes to high energy (§§17--18), with a unique-vacuum realisation of a selected macroscopic limit of a second observable (§20). These are statements of the workbench and are not verified here. The aim of the direction, as the author describes it, is an object that has a gap at every finite regulator but whose continuum limit has none, because in the limit the spectrum becomes continuous. The $S^6$ period data supply the background that drives the sequences, and the author proposes extending the construction to higher-dimensional period data related to the Leech lattice (not verified here).
My assessment: the mechanism the spatial line exhibits, finite gaps that close along a joint limit while spectral weight escapes, is the phenomenon that separates a lattice gap from a continuum gap, so I think the direction is well posed. Its next requirement is a sequence of this kind whose limit satisfies the reconstruction axioms under which the continuum problem is posed (\cite{JaffeWitten}, §4), since the present sequences are selected diagonals (Question~\ref{q:axioms}).

\paragraph{Toward weak coupling.}
Every result of the workbench holds at a fixed lattice spacing and strong bare coupling, and Proposition~\ref{prop:path} says how far along the standing path the estimates reach. The spatial line's conclusion names the interacting local continuum and a nonzero low-energy spectral weight as its target (§28).
My assessment: I think the most transferable ingredients of the present proof are Lemma~\ref{lem:casimir}, which is purely combinatorial and holds for every coupling, and the ground-state transform of Lemma~\ref{lem:gst}, which is exact; what is specific to strong coupling is the convergence of the vacuum expansion in local norms. A route toward weak coupling would have to replace that expansion by a multiscale control of the vacuum, and the blocking maps of the quantum line are a natural starting point for that (Question~\ref{q:multiscale}).

\paragraph{Lower thresholds.}
The workbench names the same-support linearized action $J_5$ with a preconditioned residual as its next candidate (§H7). For certificates built only from majorant series with nonnegative coefficients, Pringsheim's theorem suggests that re-centring cannot pass $\alpha_N$, and the workbench's calculation shows this for $N=1$. Whether signed or same-support information can pass it is Question~\ref{q:recentre}.

\paragraph{Independent confirmation of the inputs.}
An implementation independent of the workbench's programs would remove the condition in Theorem~\ref{thm:ym01}. The referee pass of version~1 computed the order-3 norms numerically with its own implementation (Clebsch--Gordan contraction over the $612$ anchored labels, in $22$ symmetry classes): $m(v_3)=1203.1424\ldots$ and $t(v_3)=127.2886\ldots$, below the workbench's $1611.59$ and $180.35$, with the record's energy coefficient $e_4=5M/216-2J/1053$ reproduced exactly. If confirmed in exact arithmetic, orders 1--3 give $\Delta_L\ge\kappa d_3(\xi)$ for $g^2\ge3.7154$ (exactly $3.71533\ldots$), with $d_3(1/64)=1.902$, recomputed here from the referee's values rounded up to $(m_3,t_3)=(1203.1425,127.2887)$. The computation of the norms is reported by the referee and was not repeated here (Question~\ref{q:inputs}).

\paragraph{The first-order loss.}
The true gap has no first-order term (Proposition~\ref{prop:gapupper}(c)), while the lower bounds lose one, because the spectral return bounds $2\Gamma(v,h')$ by $4t(v)\|Kh'\|_X$.
My assessment: I think a second-order lower bound $\Delta_L\ge\kappa(3-c\xi^2)$ is within reach of the same method, because the loss comes from one estimate that ignores the vanishing of the first-order perturbation on the plaquette eigenspace (Question~\ref{q:secondorder}).

\paragraph{Clustering.}
Suppose, as §H34 asserts, that the coefficient of $\xi^n$ in $\widehat C_{pq}$ vanishes unless the face distance of $p$ and $q$ is at most $n$. Then the bounds of Section~\ref{sec:heat} give $|\widehat C_{pq}(\tau;\xi)|\le\frac{2829}{13}e^{-13\tau/8}(55|\xi|)^{2\lceil d(p,q)/2\rceil}/(1-(55|\xi|)^2)$, uniformly in $L$. A written proof of that locality statement is open (Question~\ref{q:locality}).

\paragraph{The bridges.}
The typed maps of Section~\ref{sec:side} bring the Jacobian flow, Navier--Stokes solutions and the $S^6$ period data into the Hamiltonian setting.
My assessment: of these, I think the $S^6$ map is the most economical, since it needs one positive number $D$, and the Navier--Stokes map the most developed, since its rates are now derived conditionally \cite{NSPaper}. For the Jacobian map, the result that its trial states track the free two-gluon threshold identifies what the map produces; which properties of $F$ would move the quotients outside the window $1\pm1/(10j)$ is Question~\ref{q:jacwindow}.
